% ==============================================
% PHẦN 3 (Slide 21--30): Hàm mất mát, Metric đánh giá,
% Gradient và Optimizer trong huấn luyện 3DGS
% ==============================================

% --- Slide 21 ---
\begin{frame}{Hàm mất mát: kết hợp L1 và D-SSIM}
    \begin{columns}[T]
        \begin{column}{0.5\textwidth}
            \begin{itemize}
                \item Loss huấn luyện kết hợp hai thành phần bổ trợ nhau:
                \begin{itemize}
                    \item \textbf{L1}: sai khác trung bình theo từng pixel.
                    \item \textbf{D-SSIM}: $1-\text{SSIM}$, phản ánh sai lệch \textbf{cấu trúc}.
                \end{itemize}
                \item Công thức tổng:
                \[
                    \mathcal{L} = (1-\lambda)\,\mathcal{L}_1 + \lambda\,(1-\text{SSIM})
                \]
                \item Hệ số $\lambda = $ \texttt{lambda\_dssim}:
                \begin{itemize}
                    \item Mặc định: $0.2$
                    \item \texttt{fastgs-lite}: $0.25$ (nhấn mạnh cấu trúc hơn)
                \end{itemize}
            \end{itemize}
        \end{column}
        \begin{column}{0.5\textwidth}
            \centering
            \includegraphics[width=0.95\linewidth]{ch5_loss_mix.png}
            \captionof{figure}{Tỉ trọng L1 và D-SSIM trong loss}
        \end{column}
    \end{columns}
\end{frame}

% --- Slide 22 ---
\begin{frame}{SSIM: so sánh theo cửa sổ trượt}
    \begin{columns}[T]
        \begin{column}{0.5\textwidth}
            \begin{itemize}
                \item SSIM không so sánh pixel đơn lẻ mà trượt một \textbf{cửa sổ nhỏ} (thường $11\times 11$) trên ảnh render và ảnh gốc.
                \item Trong mỗi cửa sổ, kết hợp 3 yếu tố:
                \begin{itemize}
                    \item Độ sáng (luminance)
                    \item Độ tương phản (contrast)
                    \item Cấu trúc (structure)
                \end{itemize}
                \item Nhạy với sai lệch \textbf{cấu trúc cục bộ} hơn hẳn L1/L2 thuần pixel.
                \item Do đó bổ sung tốt cho L1 trong hàm loss.
            \end{itemize}
        \end{column}
        \begin{column}{0.5\textwidth}
            \centering
            \includegraphics[width=0.95\linewidth]{ch5_ssim_window.png}
            \captionof{figure}{Cửa sổ trượt khi tính SSIM}
        \end{column}
    \end{columns}
\end{frame}

% --- Slide 23 ---
\begin{frame}{PSNR: metric per-pixel}
    \begin{columns}[T]
        \begin{column}{0.5\textwidth}
            \begin{itemize}
                \item PSNR (Peak Signal-to-Noise Ratio) đo theo dB, dựa trên MSE giữa ảnh render và ảnh gốc:
                \[
                    \text{PSNR} = 10 \cdot \log_{10}\!\left(\frac{\text{MAX}^2}{\text{MSE}}\right)
                \]
                \item PSNR càng cao $\Rightarrow$ sai số pixel càng nhỏ.
                \item Nhược điểm: là metric \textbf{per-pixel thuần}, không phản ánh tốt chất lượng \textbf{cảm nhận thị giác} như SSIM hay LPIPS.
            \end{itemize}
        \end{column}
        \begin{column}{0.5\textwidth}
            \centering
            \includegraphics[width=0.95\linewidth]{ch5_psnr.png}
            \captionof{figure}{Quan hệ MSE và PSNR (dB)}
        \end{column}
    \end{columns}
\end{frame}

% --- Slide 24 ---
\begin{frame}{Điểm số tổng hợp (Score) đánh giá benchmark}
    \begin{columns}[T]
        \begin{column}{0.5\textwidth}
            \footnotesize
            \begin{itemize}
                \item Công thức điểm dùng trong benchmark/cuộc thi, kết hợp cả 3 metric:
                \[
                    \text{Score} = 0.4(1-\text{LPIPS}) + 0.3\,\text{SSIM} + 0.3\,\text{clamp}\!\left(\frac{\text{PSNR}}{\text{PSNR}_{\max}}, 0, 1\right)
                \]
                với $\text{PSNR}_{\max} = 30$dB.
                \item Mục tiêu: cân bằng giữa
                \begin{itemize}
                    \item độ chính xác pixel (PSNR),
                    \item cấu trúc (SSIM),
                    \item cảm nhận thị giác (LPIPS).
                \end{itemize}
            \end{itemize}
        \end{column}
        \begin{column}{0.5\textwidth}
            \centering
            \includegraphics[width=0.95\linewidth]{ch5_score_decomp.png}
            \captionof{figure}{Phân rã đóng góp của Score}
        \end{column}
    \end{columns}
\end{frame}

% --- Slide 25 ---
\begin{frame}{Lan truyền gradient qua rasterizer khả vi}
    \begin{itemize}
        \item Phép render 3DGS (chiếu Gaussian + alpha blending) hoàn toàn \textbf{khả vi}, nên có thể lan truyền ngược (backprop) từ loss về từng tham số Gaussian.
        \item Chuỗi lan truyền ngược:
    \end{itemize}
    \begin{center}
        \footnotesize
        pixel $\xrightarrow{\ \partial\ }$ alpha blend $\xrightarrow{\ \partial\ }$ alpha field $\xrightarrow{\ \partial\ }$ covariance chiếu $\xrightarrow{\ \partial\ }$ $(\mu, \Sigma)$ gốc
    \end{center}
    \begin{itemize}
        \item Quy tắc chuỗi (chain rule) tổng quát cho mỗi tham số $\theta \in \{\mu, \Sigma, \alpha, \text{SH}\}$:
        \[
            \frac{\partial \mathcal{L}}{\partial \theta} = \frac{\partial \mathcal{L}}{\partial C} \cdot \frac{\partial C}{\partial \theta}
        \]
        \item Nhờ đó mọi tham số Gaussian đều được cập nhật trực tiếp từ tín hiệu loss ở không gian ảnh.
    \end{itemize}
\end{frame}

% --- Slide 26 ---
\begin{frame}{Gradient tập trung tại biên Gaussian}
    \begin{columns}[T]
        \begin{column}{0.5\textwidth}
            \begin{itemize}
                \item Gradient view-space lớn nhất xuất hiện ở vùng \textbf{biên/rìa} của Gaussian.
                \item Đây là nơi ảnh render lệch nhiều nhất so với ảnh thật (sai số cao).
                \item Ý nghĩa quan trọng: gradient tại biên chính là \textbf{tín hiệu} dùng để quyết định \textbf{densify} (thêm Gaussian mới) ở giai đoạn sau.
                \item Vùng gradient thấp $\Rightarrow$ Gaussian đã khớp tốt, không cần chia tách thêm.
            \end{itemize}
        \end{column}
        \begin{column}{0.5\textwidth}
            \centering
            \includegraphics[width=0.95\linewidth]{ch6_edge_gradient.png}
            \captionof{figure}{Gradient tập trung ở biên Gaussian}
        \end{column}
    \end{columns}
\end{frame}

% --- Slide 27 ---
\begin{frame}{Adam optimizer cho từng tham số Gaussian}
    \begin{columns}[T]
        \begin{column}{0.5\textwidth}
            \begin{itemize}
                \item Mỗi tham số Gaussian ($\mu, \Sigma, \alpha, \text{SH}$) có trạng thái Adam riêng:
                \begin{itemize}
                    \item Moment bậc 1 (trung bình gradient)
                    \item Moment bậc 2 (trung bình bình phương gradient)
                \end{itemize}
                \item Learning rate thích ứng (adaptive) theo từng chiều tham số.
                \item Quỹ đạo hội tụ \textbf{mượt hơn} so với SGD thuần, tránh dao động mạnh khi gradient nhiễu.
            \end{itemize}
        \end{column}
        \begin{column}{0.5\textwidth}
            \centering
            \includegraphics[width=0.95\linewidth]{ch6_adam_trajectory.png}
            \captionof{figure}{Quỹ đạo cập nhật của Adam}
        \end{column}
    \end{columns}
\end{frame}

% --- Slide 28 ---
\begin{frame}{Lịch suy giảm learning rate}
    \begin{columns}[T]
        \begin{column}{0.5\textwidth}
            \begin{itemize}
                \item Learning rate của từng nhóm tham số (đặc biệt là vị trí $\mu$) giảm dần theo \textbf{hàm mũ} (exponential decay).
                \item Suy giảm diễn ra qua số bước \texttt{position\_lr\_max\_steps}.
                \item Mục đích:
                \begin{itemize}
                    \item Học nhanh, bước lớn ở giai đoạn đầu.
                    \item Ổn định, tinh chỉnh nhỏ dần về cuối huấn luyện.
                \end{itemize}
            \end{itemize}
        \end{column}
        \begin{column}{0.5\textwidth}
            \centering
            \includegraphics[width=0.95\linewidth]{ch6_lr_decay.png}
            \captionof{figure}{Suy giảm learning rate theo hàm mũ}
        \end{column}
    \end{columns}
\end{frame}

% --- Slide 29 ---
\begin{frame}{Sparse Adam schedule (đòn bẩy tăng tốc \#3)}
    \begin{columns}[T]
        \begin{column}{0.5\textwidth}
            \footnotesize
            \begin{itemize}
                \item Không phải mọi tham số được cập nhật ở mọi iteration.
                \item \texttt{features\_rest} (SH bậc cao, chi phối màu view-dependent):
                \begin{itemize}
                    \item Chỉ \texttt{step} mỗi \textbf{16} iteration đến 15k.
                \end{itemize}
                \item Sau đó, mọi tham số:
                \begin{itemize}
                    \item \texttt{step} mỗi \textbf{32} iteration (15k--20k)
                    \item \texttt{step} mỗi \textbf{64} iteration (sau 20k)
                \end{itemize}
                \item Giảm mạnh chi phí \texttt{optimizer.step()}, gần như không ảnh hưởng chất lượng vì thành phần SH bậc cao thay đổi chậm.
            \end{itemize}
        \end{column}
        \begin{column}{0.5\textwidth}
            \centering
            \includegraphics[width=0.95\linewidth]{ch6_step_schedule.png}
            \captionof{figure}{Lịch step thưa của Sparse Adam}
        \end{column}
    \end{columns}
\end{frame}

% --- Slide 30 ---
\begin{frame}{Tổng kết vòng lặp huấn luyện}
    \begin{center}
        \footnotesize
        \begin{tikzpicture}[
            node distance=1.1cm and 1.6cm,
            box/.style={draw, rounded corners, fill=blue!8, align=center, minimum width=2.6cm, minimum height=1cm},
            ->, thick
        ]
            \node[box] (render) {1. Render};
            \node[box, right=of render] (loss) {2. Tính loss\\ (L1 + D-SSIM)};
            \node[box, right=of loss] (backward) {3. Backward\\ (lan gradient)};
            \node[box, below=of loss] (step) {4. optimizer.step()\\ (sparse Adam)};
            \node[box, left=of step] (dp) {5. Densify /\\ Prune (định kỳ)};

            \draw (render) -- (loss);
            \draw (loss) -- (backward);
            \draw (backward) |- (step);
            \draw (step) -- (dp);
            \draw (dp) |- (render);
        \end{tikzpicture}
    \end{center}
    \begin{itemize}
        \item Một iteration: \textbf{render} $\to$ \textbf{tính loss} $\to$ \textbf{backward} $\to$ \textbf{optimizer.step()} $\to$ (định kỳ) \textbf{densify/prune} $\to$ lặp lại.
        \item Vòng lặp này chạy hàng chục nghìn iteration cho đến khi hội tụ.
    \end{itemize}
\end{frame}
